\section[反常层与双曲局部化]{反常层与双曲局部化（1 月 25 日）}

\srcpage{118}

我们的目标是得到几何 Satake 定理，它（粗略地）断言
\[
(\Perv_{L^+G}(\Gr_G,\Zl),\star)\cong(\Rep\widehat G,\otimes).
\]
然而今天我们将把大部分时间花在一些基础材料上。

\subsection{反常层的回顾}

我们先在代数几何的框架下回顾经典理论。设 $X$ 是代数闭域 $k$ 上的分离、有限型概形。

设 $\Lambda$ 是被某个 $n\in k^\times$ 杀灭的 Noether 环。我们也可以考虑 $\Lambda=\Ql$，但这里想集中于挠的情形。

我们有 étale 层的导出范畴 $D_\et(X,\Lambda)=D(X_\et,\Lambda)$。它紧生成，其紧对象恰为
$D^b_{c,\mathrm{ftor}}(X_\et,\Lambda)$：上同调层可构造（constructible），并在 $\Lambda$ 上具有有限 Tor 维数的有界复形。

\begin{remark}
我们有 $D^b_{c,\mathrm{ftor}}(X_\et,\Lambda)\subset D^b_c(X_\et,\Lambda)$，后者是上同调层可构造的一切有界复形构成的子范畴。
\end{remark}

\begin{definition}
\begin{enumerate}[label=(\arabic*)]
\item 满子范畴 ${}^pD^{\le0}_\et(X,\Lambda)\subseteq D_\et(X,\Lambda)$ 由这样的 $A\in D_\et(X,\Lambda)$ 组成：对一切几何点 $\overline{x}\to X$，有 $A_{\overline{x}}\in D^{\le -d(\overline{x})}(\Lambda)$，其中 $d(\overline{x})=\dim\{\overline{x}\}=\trdeg k(\overline{x})/k$。
\item 满子范畴 ${}^pD^{\le n}_\et(X,\Lambda)\subseteq D_\et(X,\Lambda)$ 指 ${}^pD^{\le0}_\et(X,\Lambda)[-n]$。
\item 满子范畴 ${}^pD^{\ge0}_\et(X,\Lambda)\subseteq D_\et(X,\Lambda)$ 是 ${}^pD^{\ge-1}_\et(X,\Lambda)$ 的右正交补，即 $B\in{}^pD^{\ge0}$ 当且仅当对一切 $A\in{}^pD^{\le-1}$ 有 $\Hom(A,B)=0$。
\item 满子范畴 ${}^pD^{\ge n}_\et(X,\Lambda)\subseteq D_\et(X,\Lambda)$ 指 ${}^pD^{\ge0}_\et(X,\Lambda)[-n]$。
\end{enumerate}
\end{definition}

\begin{theorem}
\begin{enumerate}[label=(\arabic*)]
\item 二元组 $({}^pD^{\le0}_\et,{}^pD^{\ge0}_\et)$ 定义 $D_\et(X,\Lambda)$ 上的一个 t-结构。特别地，存在截断函子
\[
{}^p\tau^{\ge0}\colon D_\et(X,\Lambda)\to{}^pD^{\le0}_\et(X,\Lambda),\qquad
{}^p\tau^{\le0}\colon D_\et(X,\Lambda)\to{}^pD^{\ge0}_\et(X,\Lambda),
\]
它们是包含函子的左/右伴随，并且
\[
{}^p\tau^{\le0}A\to A\to{}^p\tau^{\ge1}A
\]
是正合三角形。\yizhu{原文在定义 22.2(3) 与此处的 $\le$、$\ge$ 上标上似有笔误：截断函子的目标写作 ${}^pD^{\le0}$、${}^pD^{\ge0}$，两个上标恰与通常约定相反；定义 22.2(3) 中「右正交补」的对象印作 ${}^pD^{\ge-1}$，紧随其后的即字句却用 ${}^pD^{\le-1}$。译文均照录原文，读者可按通常约定理解。}{A（原文 §22.1 p.118）}
\item $A\in D_\et(X,\Lambda)$ 落在 ${}^pD^{\ge0}_\et$ 中，当且仅当对一切几何点 $i_{\overline{x}}\colon\overline{x}\to X$，有 $Ri^!_{\overline{x}}A\in D^{\ge -d(\overline{x})}(\Lambda)$。按定义，若把 $i_{\overline{x}}$ 分解为
\[
\overline{x}\xrightarrow{j_{\overline{x}}}\{\overline{x}\}\xrightarrow{i}X,
\]
则 $Ri^!_{\overline{x}}A:=j^*_{\overline{x}}Ri^!A$。
\item 它诱导 $D^b_c(X,\Lambda)$ 上的一个 t-结构（等价地，截断函子 ${}^p\tau^{\ge0}$、${}^p\tau^{\le0}$ 保持这个子范畴）。
\end{enumerate}
\end{theorem}

\begin{warning}
截断函子 ${}^p\tau^{\ge0}$、${}^p\tau^{\le0}$ 一般并不保持 $D^b_{c,\mathrm{ftor}}(X,\Lambda)$。当 $X=\Spec k$ 时，这种微妙之处已经出现：完美 $\Lambda$-复形的截断不必是完美的。不过，若 $\Lambda$ 是正则环，则没有问题。
\end{warning}

\begin{definition}
反常层范畴是上述 t-结构的心（heart），即
\[
\Perv(X,\Lambda):=({}^pD^{\le0}_\et\cap{}^pD^{\ge0}_\et).
\]
\srcpage{119}
它是一个 abelian 范畴（这是 t-结构的一般事实）。
\end{definition}

\begin{example}
若 $i\colon\Spec k\hookrightarrow X$，则 $i_*\Lambda$ 是反常的。
\end{example}

\begin{example}
若 $X$ 光滑、维数为 $d$，则 $\Lambda[d]$ 是反常的。这最终来自纯度（purity）。
\end{example}

\begin{theorem}
若 $\Lambda=\overline{\mathbb{F}}_\ell$，则 $\Perv(X,\Lambda)\cap D^b_c(X,\Lambda)$ 是阿廷范畴，即每个对象都有有限长度。不可约对象与如下二元组 $(Z,\rho)$ 一一对应：$Z$ 是 $X$ 的闭不可约子集，$\rho$ 是 $k(Z)$ 的绝对 Galois 群在 $\overline{\mathbb{F}}_\ell$-向量空间上的不可约表示。
\end{theorem}

\emph{证明概要。} 给定 $Z$ 与 $\rho$，我们得到稠密开子集 $U\subset Z$ 及浸入 $j\colon U\hookrightarrow Z$，以及 $U$ 上不可约的 $\overline{\mathbb{F}}_\ell$-局部系统 $L$。再作进一步限制，可以假设 $U$ 是光滑的。

记 $d_Z:=\dim Z$。由例 22.7，$L[d_Z]\in\Perv(U,\overline{\mathbb{F}}_\ell)$。容易验证 $j_!L[d_Z]\in{}^pD^{\le0}(Z,\overline{\mathbb{F}}_\ell)$ 且 $Rj_*L[d_Z]\in{}^pD^{\ge0}(Z,\overline{\mathbb{F}}_\ell)$。我们定义
\[
{}^pj_!L[d_Z]:={}^p\tau^{\ge0}(j_!L[d_Z]),\qquad
{}^pRj_*L[d_Z]:={}^p\tau^{\le0}(Rj_*L[d_Z]).
\]
然后定义相交复形 $\mathrm{IC}(Z,L)$ 为 ${}^pj_!L[d_Z]\to{}^pRj_*L[d_Z]$ 的像。\yizhu{原文此处把 ${}^pj_!$ 与 ${}^pRj_*$ 分别定义为 ${}^p\tau^{\ge0}(j_!L[d_Z])$ 与 ${}^p\tau^{\le0}(Rj_*L[d_Z])$。按上一段所述，$j_!L[d_Z]$ 落在 ${}^pD^{\le0}$ 中、$Rj_*L[d_Z]$ 落在 ${}^pD^{\ge0}$ 中，故通常的两个截断上标应分别取 $\le0$ 与 $\ge0$；此处疑与定理 22.3 同属笔误。译文照录原文。}{A（原文 §22.1 p.119）}

于是得到反常层 $i_*\mathrm{IC}(Z,L)\in\Perv(X,\overline{\mathbb{F}}_\ell)$。这些正是不可约对象。关键在于这一步是良定义的，且独立于 $U$ 的选取。（可构造层没有这种独立性，这正是可构造层不构成阿廷范畴的原因。）\hfill$\square$

\subsection{相对反常层}

设 $f\colon X\to S$ 是分离、有限型的态射，$S$ 是任意概形。目标：定义「相对于 $S$ 的反常性」这一概念。

\begin{definition}
\begin{enumerate}[label=(\arabic*)]
\item ${}^{p/S}D^{\le0}_\et(X,\Lambda)\subset D_\et(X,\Lambda)$ 是由满足下述条件的 $A\in D_\et(X,\Lambda)$ 构成的满子范畴：对一切几何点 $\overline{s}\to S$，有 $A|_{X_{\overline{s}}}\in{}^pD^{\le0}(X_{\overline{s}},\Lambda)$。等价地，对一切位于几何点 $\overline{s}\to S$ 上方的几何点 $\overline{x}\to X$，有 $A_{\overline{x}}\in D^{\le -d(\overline{x}/\overline{s})}(\Lambda)$。
\item ${}^{p/S}D^{\ge0}(X,\Lambda)$ 是 ${}^{p/S}D^{\le-1}$ 的右正交补。
\end{enumerate}
\end{definition}

\begin{theorem}[Hansen-S]
\begin{enumerate}[label=(\arabic*)]
\item 上述定义给出 $D_\et(X,\Lambda)$ 上的一个 t-结构。\yizhu{「Hansen-S」是原文对「Hansen 与 Scholze」的缩写，指他们关于相对反常层的结果。}{A（原文 §22.2 p.119）}
\item $A\in D_\et(X,\Lambda)$ 落在 ${}^{p/S}D^{\ge0}(X,\Lambda)$ 中，当且仅当对一切几何点 $\overline{s}\to S$，$*$-限制 $A|_{X_{\overline{s}}}\in{}^{p/S}D^{\ge0}(X_{\overline{s}},\Lambda)$。
\item 它诱导 $D^b_c(X,\Lambda)$ 上的一个 t-结构。
\end{enumerate}
\end{theorem}

\begin{remark}
当 $S$ 是离散赋值环（DVR）时，该定理的内容本质上等价于 Gabber 关于反常层在邻近轮（nearby cycles）下具有良好性质的结论。
\end{remark}

\begin{corollary}
沿 $S'\to S$ 的拉回给出 t-正合函子。更确切地说，给定笛卡尔图表
\[
\begin{tikzcd}
X' \arrow[r] \arrow[d] & X \arrow[d] \\
S' \arrow[r] & S
\end{tikzcd}
\]
$*$-拉回给出
\[
{}^{p/S}D^{\le0}(X,\Lambda)\to{}^{p/S'}D^{\le0}(X',\Lambda),
\]
\srcpage{120}
\begin{center}
\figP{120}
\end{center}
以及
\[
{}^{p/S}D^{\ge0}(X,\Lambda)\to{}^{p/S'}D^{\ge0}(X',\Lambda).
\]
\end{corollary}

\begin{corollary}
存在「$X/S$ 上的反常层族」这一概念，即
\[
\Perv(X/S,\Lambda):={}^{p/S}D^{\le0}(X,\Lambda)\cap{}^{p/S}D^{\ge0}(X,\Lambda).
\]
\end{corollary}

\subsection{$p$ 进几何中的反常层}

$p$ 进几何中的情形是微妙的。一个直接的问题是：如何定义一个点的维数。

\begin{warning}
我不知道如何定义 $\mathbb{B}^2_{C_p}$ 的一个点的「正确」维数。
\end{warning}

\begin{example}
考虑 $|\mathbb{B}^{\mathrm{ad}}_{C_p}|$。
\begin{center}
\figP{120}
\end{center}

经典点的维数为 $0$。其余秩 $1$ 的点都应当是 $1$ 维的。

那么秩 $2$ 的点呢？依视角不同，它们看起来或是 $0$ 维或是 $1$ 维。这两种取法在 Verdier 对偶下互换。
\end{example}

\begin{example}
考虑 $|\mathbb{B}^2_{C_p}|$。秩 $1$ 的点没有分类。「拓扑超越次数（topological transcendence degree）」的行为很怪异，例如存在塔 $C_p\subset K_1\subset K_2$，其中 $K_1/C_p$ 的拓扑超越次数为 $1$，$K_2/K_1$ 的拓扑超越次数为 $1$，而 $K_2/C_p$ 的拓扑超越次数也是 $1$ [Tem]。若某点的剩余域是 $K_2$，它该算 $1$ 维还是 $2$ 维？
\end{example}

小结：在 $p$ 进几何中，似乎没有希望建立一个完全一般的反常层理论。

另一方面，我们并不需要完全一般的理论。我们只需要 $\mathrm{Hk}_G\to\Div^1$ 的「相对反常层」理论。因此只需定义 $\mathrm{Hk}_G\times_{\Div^1}\Spa C$ 的点的维数。为此我们使用 Cartan 分层，并手工定义维数。我们有 $\mathrm{Hk}_G=L^+G\backslash\Gr_G$，而 $\Gr_G$ 上的 $L^+G$-轨道诱导出分层
\[
\Gr_G=\bigsqcup_{\mu\in X_*(T)_+}\Gr_{G,\mu},
\]
且 $\dim\Gr_{G,\mu}=\langle2\rho,\mu\rangle$。

\subsection{双曲局部化}

\srcpage{121}
\begin{center}
\figP{121A}
\end{center}
\begin{center}
\figP{121B}
\end{center}

我们先讲概形情形的经典背景。设 $k$ 是代数闭域，$X/k$ 是带 $\mathbb{G}_m$-作用的紧合概形。

不动点 $X^0=X^{\mathbb{G}_m}$ 构成 $X$ 的闭子集。写
\[
X^0=\bigsqcup_{i=1}^{m}X^0_i,
\]
其中每个 $X^0_i$ 在 $X^0$ 中既开又闭\footnote{通常直接取 $X^0_i$ 为 $X^0$ 的连通分支即可。然而用这种更灵活的方式陈述，能更好地适应 $p$ 进几何。}。据此，我们可以定义两个 $\mathbb{G}_m$-稳定的分层：
\begin{enumerate}[label=(\arabic*)]
\item 分层 $X=\bigcup_{i=1}^m X^+_i$，其中 $X^+_i$ 由这样的点 $x$ 组成：「$\lim_{t\to0}t\cdot x$ 存在，且落在 $X^0_i$ 中」。
\item 分层 $X=\bigcup_{i=1}^m X^-_i$，其中 $X^-_i$ 由这样的点 $x$ 组成：「$\lim_{t\to\infty}t\cdot x$ 存在，且落在 $X^0_i$ 中」。
\end{enumerate}

更确切地说，设 $(\mathbb{A}^1)^+\hookrightarrow\mathbb{P}^1$ 是 $\infty$ 的补集，$(\mathbb{A}^1)^-\hookrightarrow\mathbb{P}^1$ 是 $0$ 的补集。$X^+_i$ 上的 $\mathbb{G}_m$-作用延拓为
\[
(\mathbb{A}^1)^+\times X^+_i\to X^+_i,
\]
把 $0\times X^+_i$ 送到 $X^0_i$；而 $X^-_i$ 上的 $\mathbb{G}_m$-作用延拓为
\[
(\mathbb{A}^1)^-\times X^-_i\to X^-_i,
\]
把 $0\times X^-_i$ 送到 $X^0_i$。

我们写 $X^+:=\bigsqcup X^+_i$，$X^-:=\bigcup X^-_i$。

\begin{example}
考虑 $\mathbb{G}_m$ 按通常方式作用在 $X=\mathbb{P}^1$ 上。
\begin{center}
\figP{121A}
\end{center}
\begin{center}
\figP{121B}
\end{center}

我们取分解
\[
X^0=\{0,\infty\}=\{0\}\sqcup\{\infty\},
\]
此时
\[
X^+=\mathbb{A}^1\sqcup\{\infty\},\qquad X^-=\{0\}\sqcup(\mathbb{A}^1)^-.
\]
\end{example}

\begin{exercise}
考虑 $\mathbb{G}_m$ 按 $t\cdot(a_1,a_2)=(t^{-1}a_1,ta_2)$ 作用在 $\mathbb{P}^1\times\mathbb{P}^1$ 上。求出
\[
(\mathbb{P}^1\times\mathbb{P}^1)^0=(0,0)\sqcup(0,\infty)\sqcup(\infty,0)\sqcup(\infty,\infty)
\]
所对应的分层。
\end{exercise}

双曲局部化的目标，是用 $X^0$ 处的局部信息来描述 $X$ 上 $\mathbb{G}_m$-等变层的上同调。

\srcpage{122}
\begin{center}
\figP{122}
\end{center}

\begin{theorem}[Braden]
存在一个函子
\[
L\colon D_\et(X/\mathbb{G}_m,\Lambda)\to D_\et(X^0,\Lambda),
\]
使得 $R\Gamma(X,A)$ 带有滤过，其伴随分次为 $R\Gamma(X^0,L(A))$。

事实上，$L$ 有四个显式描述。考虑笛卡尔图表
\[
\begin{tikzcd}[column sep=6.5em, row sep=3.4em]
X^0_i \arrow[r, "i^-"'] \arrow[d, "i^+"'] & X^-_i \arrow[l, "p^-", bend right=45] \arrow[d, "q^-"] \\
X^+_i \arrow[u, "p^+", bend left=45] \arrow[r, "q^+"'] & X
\end{tikzcd}
\]
函子 $L$ 有如下四个自然同构的描述\yizhu{原文此式中 $!$ 的排法不一致：$R(p^-)_!$、$R(p^+)_*$ 用下标，而 $R(i^-)^!$、$R(q^+)^!$ 印成上标；按 Braden 的记法，后者与 $(q^-)^*$、$(i^+)^*$ 对偶，仍应是下叹号 $!$。译文照录原文的上下标。又，例 22.20 末项印作 $(q^-)^+A$，其中 $+$ 疑为 $*$ 之误。}{A（原文 §22.4 p.122）}：
\[
R(p^-)_!(q^-)^*\xleftarrow{\ \sim\ }R(i^-)^!(q^-)^*\xleftarrow{\ \sim\ }(i^+)^*R(q^+)^!\xleftarrow{\ \sim\ }R(p^+)_*R(q^+)^!.
\]
\end{theorem}

\begin{example}
考虑 $\mathbb{A}^2$ 上带双曲 $\mathbb{G}_m$-作用的 $\mathbb{G}_m$-等变层 $A$，其中 $t\cdot(a_1,a_2)=(t^{-1}a_1,ta_2)$。
\begin{center}
\figP{122}
\end{center}

此时有笛卡尔图表
\[
\begin{tikzcd}
0 \arrow[r, "i^-"] \arrow[d, "i^+"'] & \mathbb{A}^1\times\{0\} \arrow[d, "q^-"] \\
0\times\mathbb{A}^1 \arrow[r, "q^+"'] & \mathbb{A}^2
\end{tikzcd}
\]
于是双曲局部化函子 $D(X/\mathbb{G}_m,\Lambda)\to D(\Lambda)$ 可以取为
\[
(i^+)^*R(q^+)^!A\xrightarrow{\ \sim\ }R(i^-)^!(q^-)^+A.
\]
\end{example}

\begin{example}
考虑 $X=\mathbb{P}^1$ 带通常的 $\mathbb{G}_m$-作用，并取 $A=\Lambda$。则 $R\Gamma(\mathbb{P}^1,\Lambda)=\Lambda[0]\oplus\Lambda[-2]$。此时
\[
L(A)_{\{\infty\}}=R\Gamma_c(\mathbb{A}^1,\Lambda)=\Lambda[-2],
\]
它同时也与 $R\Gamma_{\{\infty\}}(\mathbb{A}^1,\Lambda)=\Lambda[-2]$ 等同。另一方面，
\[
L(A)_{\{0\}}=R\Gamma_c(\{0\},\Lambda)=\Lambda[0].
\]
所以我们在这一情形显式验证了 Braden 的定理。
\end{example}

\begin{example}
设 $X$ 是旗簇（flag variety）$G/P$。它带有 $G$ 的作用，而我们可以通过支配正则余特征把该作用拉回到 $\mathbb{G}_m$。此时 $X^{\mathbb{G}_m}=X^T=W/W_P$。Braden 的定理 \srcpage{123}推出
\[
R\Gamma(X,\Lambda)=\bigoplus_{w\in\underbrace{W^P}_{W/W_P}}\Lambda[-2\ell(w)].
\]
\end{example}

我们将把它用于 $X=\Gr_{G,\le\mu}\subset\Gr_G$ 与 $\mathbb{G}_m\subset L^+G$，以理解 $\Gr_G$ 上 $L^+G$-等变反常层的上同调。

\begin{remark}
$L$ 有多种描述，这非常有用。特别地，由于它既有用左伴随给出的描述，又有用右伴随给出的描述，它与一切极限和余极限交换。这里也有相对版本。
\end{remark}

\subsection{菱形上的双曲局部化}

设定：设 $f\colon X\to S$ 是小 v-叠之间的紧合映射，可由空间菱形表示，且 $\dim.\mathrm{tr.g.}\,f<\infty$。设 $X$ 带有 $\mathbb{G}_m$-作用\footnote{adic 的 $\mathbb{G}_m$ 的函数点由 $(R,R^+)\mapsto R^\times$ 给出。}，且 $f$ 关于该作用是等变的（$S$ 上取平凡作用）。

我们假设有两个如上所述的 $\mathbb{G}_m$-等变分层 $X=\bigcup X^+_i$ 与 $X=\bigcup X^-_i$。\footnote{这样的分层未必自动存在，而在正规概形的情形下它是自动的。}

\begin{theorem}
在此情形，对一切 $A\in D_\et(X/\mathbb{G}_m,\Lambda)$，映射
\[
R(p^-)_!(q^-)^*\xleftarrow{\ \sim\ }R(i^-)^!(q^-)^*\xleftarrow{\ \sim\ }(i^+)^*R(q^+)^!\xleftarrow{\ \sim\ }R(p^+)_*R(q^+)^!
\]
都是同构，从而定义「双曲局部化函子」
\[
L_{X/S}\colon D_\et(X/\mathbb{G}_m,\Lambda)\to D_\et(X^0,\Lambda).
\]
函子 $L_{X/S}$ 与一切（余）极限（在 $\infty$-范畴的意义下）以及一切基变换 $S'\to S$ 交换；并且对下面的图表
\[
\begin{tikzcd}
X \arrow[r, "f"] & S \\
X^0 \arrow[u] \arrow[ur, "f^0"'] &
\end{tikzcd}
\]
$Rf_*$ 上存在滤过，其伴随分次为 $Rf^0_*L_{X/S}$。
\end{theorem}

\emph{证明。} 一切都由下述几何原则推出：若 $Y$ 是有限维的局部空间菱形，带有 $\mathbb{G}_m$-作用，在 $S$ 上部分紧合，且 $[Y/\mathbb{G}_m]$ 在 $S$ 上 qcqs，则 $Y$ 有两个端点；于是由第 19 讲中的消没定理，对一切 $A\in D_\et([Y/\mathbb{G}_m],\Lambda)$ 有 $R\Gamma_{\partial c}(Y,A)=0$。这并不难证：可以化归到 $\mathbb{G}_m$ 的情形，而该情形可以显式算出。\hfill$\square$

\begin{example}
考虑 $\mathbb{G}_m$ 按通常方式作用在 $\mathbb{P}^1$ 上。$R\Gamma_c(\mathbb{A}^1,A)$ 与 $R\Gamma_{\{0\}}(\mathbb{A}^1,A)$ 之差是 $R\Gamma_{\partial c}(\mathbb{G}_m,A)$。
\end{example}

\begin{remark}
我们这里大量使用了解析世界，在其中 $\mathbb{G}_m$ 不是拟紧的，并且有两个端点。
\end{remark}
